\documentclass[10pt,twocolumn]{article}
\usepackage[margin=0.75in]{geometry}
\usepackage{booktabs}
\usepackage{graphicx}
\usepackage{hyperref}
\usepackage{xcolor}
\usepackage{enumitem}
\newcommand{\result}[1]{\textcolor{red}{[RESULT: #1]}}
\newcommand{\figuretodo}[1]{\begin{center}\fbox{\parbox{0.92\linewidth}{\centering\vspace{0.35in}#1\vspace{0.35in}}}\end{center}}

\title{From Commits to Events: Visualizing Long-Horizon Agentic Software Evolution}
\author{Anonymous Authors}
\date{}

\begin{document}
\maketitle

\begin{abstract}
Coding agents increasingly modify repositories across many sessions and days,
but their development process disappears between two incompatible records:
native transcripts contain fine-grained actions that are too large to inspect,
whereas version control retains durable snapshots but discards exploration and
failed work. We present an event-resolved representation and coordinated
visualization suite that joins agent actions with Git evolution and current
line survival. The suite adapts established software-evolution views---line
maps, evolution matrices, repository maps, animated histories, survival
strata, forensic coupling, and ownership storylines---to multi-agent,
long-horizon histories with stable layout and shared time navigation. We
evaluate what process evidence commit-only analysis loses, which behavioral
patterns event granularity reveals, whether the views improve realistic review
tasks, and how the approach scales from days to months. \result{Insert the
final dataset sizes, join quality, task-study outcomes, and performance.}
\end{abstract}

\section{Introduction}
Software evolution tools traditionally observe commits and attribute them to
people. This abstraction was practical when a person retained the reasons
behind a sequence of commits and code review could examine the resulting diff.
Long-horizon coding agents break both assumptions. They can explore and modify
a repository across many sessions, yet no persistent participant necessarily
remembers why the resulting code has its present form. A commit records the
surviving edit, not the files read before it, abandoned alternatives,
verification order, repeated tool failures, or work that never reached a
commit. Native agent logs contain many of these events, but weeks of logs are
not a usable review interface.

The obvious response is to replay a transcript or place events on a timeline.
That response is incomplete. A transcript does not establish which changes
became durable, while an unconstrained replay remains linear and expensive to
inspect. Our central thesis is that agent events and repository history are
complementary evidence: the former preserves process and the latter preserves
outcome. Joining them permits questions that neither source can answer alone,
such as whether a repeatedly edited file was eventually committed, whether a
test followed its last edit, and whether agent-associated code survived.

We build this join into AgentSight's vendor-neutral session representation and
project it into coordinated views derived from three decades of software
evolution visualization. The goal is not a gallery of visual metaphors. Every
view supports a decision: locate old or unreviewed lines, identify bursty or
decaying files, inspect exploration order, measure survival, find hidden
coupling, or understand how agents and people divided the repository. Stable
spatial layouts, semantic zoom, brushing, and a shared playback cursor allow
the same evidence to be inspected at day, session, and event scale.

This paper makes three contributions:
\begin{enumerate}[leftmargin=*]
  \item an event-resolved longitudinal representation that explicitly joins
  agent process evidence, Git evolution, and current line survival without
  conflating touch, commit, and authorship;
  \item a coordinated visualization design that adapts seven established
  software-evolution traditions to long-horizon, multi-agent development; and
  \item an evaluation of information recoverability, behavioral structure,
  review utility, and interactive scalability on real development histories.
\end{enumerate}

\section{Background and Motivation}
\subsection{Why commits are insufficient}
A commit is an intentionally packaged state transition. Its compression is a
feature for integration and a limitation for process reconstruction. Multiple
reads, unsuccessful edits, test ordering, and uncommitted work can collapse
into the same diff or disappear completely. The paper will quantify this loss
rather than assuming that finer granularity is automatically useful.

\subsection{Why event logs are insufficient}
Native coding-agent histories retain timestamps, tool inputs, model usage, and
often referenced paths. They are vendor-specific, can contain repeated or
partial messages, and do not prove a durable repository outcome. Large logs
also exceed human attention. The representation must preserve uncertainty and
support progressive aggregation.

\subsection{Software-evolution visualization traditions}
Prior work includes pixel-oriented line maps, file-by-version matrices,
spatial repository metaphors, animated commit histories, contribution-survival
flows, hotspot and temporal-coupling analysis, and ownership storylines.
\result{Insert verified primary citations and closest-work comparison after
the mandatory literature review.} These traditions supply mature visual
encodings, but usually assume commits as observations and humans as authors.

\section{Evidence Model}
The representation uses three non-interchangeable sources. Agent events record
prompts, model responses, tool actions, referenced paths, statuses, and time.
Git records commits, rename-aware path changes, additions and deletions,
authors, and parent relationships. The current tree and blame record which
paths and lines survive. Each exported observation carries source identity and
join confidence.

The join first canonicalizes repository-relative paths, follows Git renames,
and assigns events to bounded candidate commit intervals. It then distinguishes
four outcomes: observed and committed, observed but not committed, committed
without an observed event, and currently surviving. Temporal proximity is not
reported as semantic causality or cryptographic authorship. Read-before-edit
edges are ordered evidence within a prompt/session and are labeled accordingly.

\section{Visualization Design}
\figuretodo{System overview: native sessions + Git history + current blame
$\rightarrow$ longitudinal representation $\rightarrow$ coordinated views.}

\subsection{Line and matrix views}
A SeeSoft-style line map compresses current lines into colored pixels encoding
age, Git author, and agent-touch evidence. A file-by-time evolution matrix
shows size and event intensity, with candidate burst, periodicity, decay, and
short-lived patterns labeled only after quantitative definitions are fixed.

\subsection{Stable repository map and playback}
A deterministic hierarchical layout anchors paths across time. Touch events
animate over the map while a shared cursor supports play, pause, scrubbing,
speed selection, and interval brushing. Search, read, edit, verification, and
commit use distinct channels. Aggregated frames bound work at overview scale;
selection loads event detail.

\subsection{Survival, forensic, and social views}
Age strata show how much code from each period survives. Hotspots combine
change frequency with complexity proxies; Git co-change and ordered
read-before-edit networks are shown separately. Ownership and storyline views
trace people and agent families through directory clusters without presenting
inferred agent association as Git authorship.

\subsection{Longitudinal navigation}
Calendar and punch-card summaries locate active periods. Vital-sign plots show
event rate, token rate, error rate, edit/read balance, and verification gaps.
Token-flow and ghost-comparison views support cost attribution and matched-run
comparison. All views link through the same time interval and path selection.

\section{Implementation}
The exporter extends the existing Rust `agent-session` library, which already
normalizes Claude, Codex, and Gemini histories. Git data is collected through
stable command-line plumbing and streamed into a bounded intermediate format.
The browser application reuses established open-source layout, charting,
network, and time-series libraries after license and capability verification.
\result{Insert final components, versions, data volumes, aggregation strategy,
and implementation size.}

\section{Evaluation}
We organize the evaluation around four questions.

\subsection{RQ1: Recoverability}
We compare event histories with rename-aware Git changes and current survival
across real multi-day repositories. Metrics include path-normalization yield,
event-to-change match coverage, ambiguity, unmatched-event categories, and
process facts absent from commit-only histories. \result{Insert datasets,
protocol, results, and threats.}

\subsection{RQ2: Behavioral structure}
We define candidate exploration, churn, verification, coupling, and survival
patterns before inspecting final outcomes, then measure recurrence and
stability across sessions and agents. Commit-only ablations test whether event
granularity is necessary. \result{Insert definitions and final results.}

\subsection{RQ3: Review utility}
Participants answer process-review questions using a native event table,
commit-centric baselines, and the coordinated views. We measure accuracy,
completion time, confidence calibration, and interaction traces. Tasks include
locating unverified final edits, identifying abandoned work, explaining a
cross-file change, and finding high-churn code that survived. \result{Insert
approved study protocol, population, statistics, and outcomes.}

\subsection{RQ4: Long-horizon scalability}
We measure export throughput, output size, initial render time, brushing and
scrubbing latency, and memory across increasing histories. We separately
evaluate overview aggregation and detail-on-demand. \result{Insert hardware,
history sizes, thresholds, and measurements.}

\section{Discussion}
The joined evidence can preserve part of a development process, not an agent's
private reasoning or a complete causal explanation. Missing logs, clock skew,
parallel work, path aliases, and squash/rebase workflows limit attribution.
The design therefore treats mismatch as evidence rather than silently filling
gaps. If event-resolved views improve only presentation but not review tasks,
the result would bound the scientific claim even if the gallery remains useful
as an exploratory artifact.

\section{Related Work}
The final paper will compare against software evolution visualization,
repository mining and socio-technical analysis, large-trace visualization,
coding-agent observability, and trajectory evaluation. \result{Populate only
from verified primary sources and explicitly identify same-claim work.}

\section{Conclusion}
Long-horizon coding agents make software development more observable in one
sense---their native event streams are recorded---and less understandable in
another: no persistent person necessarily holds the theory behind the code.
This work tests whether joining those events with durable repository history
and projecting them through coordinated evolution views can recover a
practical review interface. \result{State final supported scope and findings.}

\end{document}
